\documentclass[
11pt, % The default document font size, options: 10pt, 11pt, 12pt
%codirector, % Uncomment to add a codirector to the title page
]{charter} 


% El títulos de la memoria, se usa en la carátula y se puede usar el cualquier lugar del documento con el comando \ttitle
\titulo{Vinculación de dispositivos IoT para monitoreo remoto de pacientes crónicos} 

% Nombre del posgrado, se usa en la carátula y se puede usar el cualquier lugar del documento con el comando \degreename
%\posgrado{Carrera de Especialización en Sistemas Embebidos} 
\posgrado{Carrera de Especialización en Internet de las Cosas} 
%\posgrado{Carrera de Especialización en Inteligencia Artificial}
%\posgrado{Maestría en Sistemas Embebidos} 
%\posgrado{Maestría en Internet de las cosas}

% Tu nombre, se puede usar el cualquier lugar del documento con el comando \authorname
% IMPORTANTE: no omitir titulaciones ni tildación en los nombres, también se recomienda escribir los nombres completos (tal cual los tienen en su documento)
\autor{Bioing. Carlos Augusto Zelayeta}

% El nombre del director y co-director, se puede usar el cualquier lugar del documento con el comando \supname y \cosupname y \pertesupname y \pertecosupname
\director{Bioing. Sergio Liberczuk}
\pertenenciaDirector{Virtual Sense S.A.} 
\codirector{} % para que aparezca en la portada se debe descomentar la opción codirector en los parámetros de documentclass
\pertenenciaCoDirector{FIUBA}

% Nombre del cliente, quien va a aprobar los resultados del proyecto, se puede usar con el comando \clientename y \empclientename
\cliente{Ing. Leandro Bergher}
\empresaCliente{Virtual Sense S.A.}
 
\fechaINICIO{29 de abril de 2025}		%Fecha de inicio de la cursada de GdP \fechaInicioName
\fechaFINALPlan{17 de junio de 2025} 	%Fecha de final de cursada de GdP
\fechaFINALTrabajo{diciembre de 2025}	%Fecha de defensa pública del trabajo final


\begin{document}

\maketitle
\thispagestyle{empty}
\pagebreak


\thispagestyle{empty}
{\setlength{\parskip}{0pt}
\tableofcontents{}
}
\pagebreak


\section*{Registros de cambios}
\label{sec:registro}


\begin{table}[ht]
\label{tab:registro}
\centering
\begin{tabularx}{\linewidth}{@{}|c|X|c|@{}}
\hline
\rowcolor[HTML]{C0C0C0} 
Revisión & \multicolumn{1}{c|}{\cellcolor[HTML]{C0C0C0}Detalles de los cambios realizados} & Fecha      \\ \hline
0      & Creación del documento                                 &\fechaInicioName \\ \hline
1      & Se completa hasta el punto 5 inclusive                & 13 de mayo de 2025 \\ \hline
2      & Se completa hasta el punto 9 inclusive \newline
		  Se implementan las correcciones de la devolución de la versión 1 
                                                  & 20 de mayo de 2025 \\ \hline
%3      & Se completa hasta el punto 12 inclusive                & {día} de {mes} de 202X \\ \hline
%4      & Se completa el plan	                                 & {día} de {mes} de 202X \\ \hline

% Si hay más correcciones pasada la versión 4 también se deben especificar acá

\end{tabularx}
\end{table}

\pagebreak


\section*{Acta de constitución del proyecto}
\label{sec:acta}

\begin{flushright}
Buenos Aires, \fechaInicioName
\end{flushright}

\vspace{2cm}

Por medio de la presente se acuerda con el \authorname\hspace{1px} que su Trabajo Final de la \degreename\hspace{1px} se titulará ``\ttitle'' y consistirá en la vinculación de instrumentos de medición de presión arterial al sistema Virtual Sense. El trabajo tendrá un presupuesto preliminar estimado de 510 horas y un costo estimado de \$ 18000000, con fecha de inicio el \fechaInicioName\hspace{1px} y fecha de presentación pública en \fechaFinalName.

Se adjunta a esta acta la planificación inicial.

\vfill

% Esta parte se construye sola con la información que hayan cargado en el preámbulo del documento y no debe modificarla
\begin{table}[ht]
\centering
\begin{tabular}{ccc}
\begin{tabular}[c]{@{}c@{}}Dr. Ing. Ariel Lutenberg \\ Director posgrado FIUBA\end{tabular} & \hspace{2cm} & \begin{tabular}[c]{@{}c@{}}\clientename \\ \empclientename \end{tabular} \vspace{2.5cm} \\ 
\multicolumn{3}{c}{\begin{tabular}[c]{@{}c@{}} \supname \\ Director del Trabajo Final\end{tabular}} \vspace{2.5cm} \\
\end{tabular}
\end{table}




\section{1. Descripción técnica-conceptual del proyecto a realizar}
\label{sec:descripcion}

La empresa Virtual Sense ofrece un servicio de monitoreo remoto de pacientes con enfermedades crónicas. Para este objetivo, la empresa posee una aplicación móvil que recolecta datos periódicamente, tales como presión arterial, saturación de oxígeno, glucemia, y otros indicadores de acuerdo a cada patología, desde dispositivos de medición vinculados a la plataforma. Estos datos son reportados a un sistema en el cual los profesionales médicos pueden acceder para realizar un seguimiento de los pacientes. Los pacientes, por su par parte, acceden también a contenido de consulta aprobado por su médico y a comunicación por chat con los profesionales.

La recolección periódica de mediciones permite el procesamiento de datos de manera continua, lo que a su vez posibilita la detección temprana y automática de arritmias cardíacas asintomáticas y otros eventos adversos. La empresa ha encontrado gran adopción por parte de médicos con pacientes que habitan lugares remotos, para los que es difícil o costoso la realización de visitas periódicas al consultorio. 

El stack de tecnología de la empresa consiste actualmente en una aplicación móvil que permite conectarse a diversos dispositivos (electrocardiógrafos, tensiómetros, etc). Esta aplicación reporta datos mediante protocolo HTTPS hacia Amazon Web Services y luego los datos son almacenados utilizando MySQL. Un panel de control web permite a los profesionales acceder a los datos de sus pacientes, así como a los pacientes acceder a sus propios datos. Este panel posee los permisos de acceso necesarios para que los usuarios solo puedan acceder a sus propios datos, y los profesionales accedan solamente a los datos de sus pacientes. En la figura \ref{fig:diagBloques} se describe brevemente el sistema implementado por la empresa para los dispositivos ya integrados.

\begin{figure}[htpb]
\centering 
\includegraphics[width=.65\textwidth]{./Figuras/Diagrama1.pdf}
\caption{Diagrama en bloques del sistema.}
\label{fig:diagBloques}
\end{figure}

La propuesta de trabajo realizada por Virtual Sense consiste en establecer una comunicación entre el sistema de la empresa y algunos dispositivos muy utilizados por los pacientes, pero para los cuales actualmente no poseen vinculación automática.  
Los tensiómetros de marca Omron son considerados dispositivos prioritarios por la empresa, dado su uso muy extendido entre los pacientes y la alta confiabilidad de las mediciones que producen. La empresa Omron ofrece varias opciones de integración de sus dispositivos, ya que pueden ser vinculados directamente mediante Bluetooth Low Energy a aplicaciones externas, o bien se puede realizar una vinculación a través de una API que permite comunicación entre la aplicación de Omron y aplicaciones externas. En base a estas opciones, se evaluará cuál es la estrategia de vinculación óptima para conectar los dispositivos a la aplicación de Virtual Sense, utilizando criterios como, por ejemplo, costo, posibilidad de mantenimiento posterior de la integración, tiempo de desarrollo y cantidad de modelos disponibles en el mercado para cada alternativa. Una vez definida la estrategia óptima para la vinculación del dispositivo, se implementará la vinculación para un dispositivo de la marca (a definir por la empresa).
La figura \ref{fig:diagAlternativas} representa las alternativas a evaluar para vincular los dispositivos:

\begin{figure}[htpb]
\centering 
\includegraphics[width=.65\textwidth]{./Figuras/Diagrama2.pdf}
\caption{Posibles alternativas para establecer la vinculación.}
\label{fig:diagAlternativas}
\end{figure}


\section{2. Identificación y análisis de los interesados}
\label{sec:interesados}


\begin{table}[ht]
%\caption{Identificación de los interesados}
%\label{tab:interesados}
\begin{tabularx}{\linewidth}{@{}|l|X|X|l|@{}}
\hline
\rowcolor[HTML]{C0C0C0} 
Rol           & Nombre y Apellido & Organización 	& Puesto 	\\ \hline
Cliente       & \clientename      &\empclientename	&    Presidente    	\\ \hline
Responsable   & \authorname       & FIUBA        	& Alumno 	\\ \hline
Orientador    & \supname	      & \pertesupname 	& Jefe científico \\ \hline
Equipo        & Lucas Ricardo Grana \newline Emilio Borovski
	&     Virtual Sense S.A.         	&  Desarrolladores      	\\ \hline
Usuario final &         -         &      -       	& Usuarios del servicio \\ \hline
\end{tabularx}
\end{table}

\begin{itemize}
\item El ingeniero Leandro Bergher es el presidente y fundador de Virtual Sense y quien determinó la necesidad de implementar la vinculación propuesta en este trabajo.
\item El bioingeniero Sergio Liberczuk conoce todos los aspectos generales del sistema y está en contacto permanente con pacientes y médicos que lo utilizan.
\item Los desarrolladores Lucas Grana y Emilio Borovski están encargados del desarrollo de la aplicación móvil y podrán asesorar en los cambios requeridos para implementar la vinculación.
\end{itemize}

\section{3. Propósito del proyecto}
\label{sec:proposito}

Lograr el envío automático de datos entre los dispositivos marca Omron y el servicio Virtual Sense. Esta funcionalidad tendrá un impacto positivo para los pacientes al evitar la carga manual de mediciones.

\section{4. Alcance del proyecto}
\label{sec:alcance}

El proyecto incluye:
\begin{itemize}
	\item Analizar las opciones de vinculación ofrecidas por Omron.
	\item Seleccionar la opción más conveniente para el sistema Virtual Sense.
	\item Implementar la vinculación seleccionada entre la aplicación Virtual Sense y un dispositivo Omron.
	\item Implementar el manejo de los datos en la nube:
	\begin{itemize}
		\item Realizar las modificaciones necesarias en la base de datos.
		\item Realizar las modificaciones en el panel de control.
	\end{itemize}
\end{itemize}

El proyecto no incluye la validación de esta comunicación en todos los dispositivos de la empresa ni el despliegue de esta funcionalidad en la versión de producción del sistema Virtual Sense.

\section{5. Supuestos del proyecto}
\label{sec:supuestos}

Para el desarrollo del presente proyecto se supone que:

\begin{itemize}
	\item Virtual Sense proveerá los dispositivos necesarios para el proyecto.
	\item Virtual Sense establecerá el contacto con Omron para la gestión de licencias en caso de ser necesario.
	\item Se dispondrá de las horas necesarias para realizar el proyecto.
\end{itemize}


\section{6. Requerimientos}
\label{sec:requerimientos}

\begin{enumerate}
	\item Requerimientos funcionales:
		\begin{enumerate}
			\item La aplicación de Virtual Sense se debe vincular con el dispositivo Omron.
			\item Los valores medidos por el dispositivo deben ser ingresados a la aplicación.
			\item El usuario debe poder consultar los valores medidos en la aplicación.
			\item Las mediciones se deben enviar a la nube, junto con la fecha y hora en la que se realizaron y la identificación del paciente.
		\end{enumerate}
	\item Requerimientos de documentación:
		\begin{enumerate}
			\item Generar las instrucciones de vinculación orientadas a los usuarios.
		\end{enumerate}
	\item Requerimiento de seguridad
		\begin{enumerate}
			\item Los datos se enviarán usando el protocolo actual de la empresa mediante HTTPS.
		\end{enumerate}
	\item Requerimientos de la interfaz en la aplicación.	
		\begin{enumerate}
			\item Se debe indicar el estado de la conexión con el dispositivo (conectado o desconectado).
			\item Se deben mostrar los valores obtenidos luego de una medición.
		\end{enumerate}
\end{enumerate}

\section{7. Historias de usuarios (\textit{Product backlog})}
\label{sec:backlog}

\begin{enumerate}

\item ``Como paciente quiero recibir una confirmación de que el tensiómetro está vinculado correctamente."

\textit{Story points}: 3 (complejidad: 1, dificultad: 1, incertidumbre: 1)

\item ``Como paciente quiero conectar mi tensiómetro con la aplicación para ver la medición en el celular."

\textit{Story points}: 13 (complejidad: 3, dificultad: 5, incertidumbre: 3)

\item ``Como paciente quiero ver mis mediciones en el panel de control para ."

\textit{Story points}: 8 (complejidad: 2, dificultad: 2, incertidumbre: 3)

\item ``Como médico quiero ver el registro de mediciones de mis pacientes para monitorear su estado."

\textit{Story points}: 8 (complejidad: 2, dificultad: 2, incertidumbre: 3)

\item ``Como médico quiero saber si una medición se cargó automáticamente o a mano para determinar la confiabilidad de cada dato"

\textit{Story points}: 8 (complejidad: 2, dificultad: 2, incertidumbre: 2)

\end{enumerate}

\section{8. Entregables principales del proyecto}
\label{sec:entregables}

Los entregables del proyecto son:

\begin{itemize}
	\item Código fuente del software de la aplicación.
	\item Código fuente del software de backend.
	\item Instrucciones de vinculación para usuarios.	
	\item Memoria del trabajo final.
\end{itemize}

\section{9. Desglose del trabajo en tareas}
\label{sec:wbs}

\begin{enumerate}
\item Investigación (80 h)
	\begin{enumerate}
	\item Lectura de documentación técnica de Omron (20 h).
	\item Aprendizaje del funcionamiento del sistema Virtual Sense (10 h).
	\item Aprendizaje de la plataforma Google Health (10 h).
	\item Aprendizaje de la plataforma Apple Health (10 h).
	\item Revisión de estrategias seguidas con tensiómetros de otras marcas ya vinculadas al sistema (20 h).
	\item Definición del modelo a vincular (10 h).
	\end{enumerate}
\item Implementación (260 h)
	\begin{enumerate}
	\item Prueba de concepto de conexión por BLE (40 h).
	\item Prueba de concepto de conexión por API (40 h).
	\item Investigar escalabilidad de las conexiones a otros dispositivos de la marca (20 h).
	\item Implementación de la vinculación elegida (40 h).
	\item Programación del envío de datos a la nube (40 h).
	\item Programación de modificaciones en el backend (40 h).
	\item Implementación de modificaciones en panel de control (40 h).
	\end{enumerate}
\item Verificación (90 h)
	\begin{enumerate}
	\item Testeo de la conexión implementada entre tensiómetro y aplicación(20 h).
	\item Testeo del envío de datos hacia la nube (10 h).
	\item Testeo de modificaciones en panel de control (10 h).
	\item Testeo de modificaciones en el backend (20 h).
	\item Testeo de conexión a servicios Google Health y Apple Health (30 h).
	\end{enumerate}	
\item Documentación (80 h)
	\begin{enumerate}
	\item Escritura de la memoria de proyecto (40 h).
	\item Escritura de instrucciones para clientes (20 h).
	\item Elaboración de la presentación (20 h).
	\end{enumerate}
\end{enumerate}

Cantidad total de horas: 510.

\section{10. Diagrama de Activity On Node}
\label{sec:AoN}

\begin{consigna}{red}
Armar el AoN a partir del WBS definido en la etapa anterior.

Una herramienta simple para desarrollar los diagramas es el Draw.io (\url{https://app.diagrams.net/}).
\href{https://app.diagrams.net}{Draw.io}


\begin{figure}[htpb]
\centering 
\includegraphics[width=.8\textwidth]{./Figuras/AoN.png}
\caption{Diagrama de \textit{Activity on Node}.}
\label{fig:AoN}
\end{figure}

Indicar claramente en qué unidades están expresados los tiempos.
De ser necesario indicar los caminos semi críticos y analizar sus tiempos mediante un cuadro.
Es recomendable usar colores y un cuadro indicativo describiendo qué representa cada color.

\end{consigna}

\section{11. Diagrama de Gantt}
\label{sec:gantt}

\begin{consigna}{red}
Existen muchos programas y recursos \textit{online} para hacer diagramas de Gantt, entre los cuales destacamos:

\begin{itemize}
\item Planner
\item GanttProject
\item Trello + \textit{plugins}. En el siguiente link hay un tutorial oficial: \\ \url{https://blog.trello.com/es/diagrama-de-gantt-de-un-proyecto}
\item Creately, herramienta online colaborativa. \\\url{https://creately.com/diagram/example/ieb3p3ml/LaTeX}
\item Se puede hacer en latex con el paquete \textit{pgfgantt}\\ \url{http://ctan.dcc.uchile.cl/graphics/pgf/contrib/pgfgantt/pgfgantt.pdf}
\end{itemize}

Pegar acá una captura de pantalla del diagrama de Gantt, cuidando que la letra sea suficientemente grande como para ser legible. 
Si el diagrama queda demasiado ancho, se puede pegar primero la ``tabla'' del Gantt y luego pegar la parte del diagrama de barras del diagrama de Gantt.

Configurar el software para que en la parte de la tabla muestre los códigos del EDT (WBS).\\
Configurar el software para que al lado de cada barra muestre el nombre de cada tarea.\\
Revisar que la fecha de finalización coincida con lo indicado en el Acta Constitutiva.

En la figura \ref{fig:gantt}, se muestra un ejemplo de diagrama de gantt realizado con el paquete de \textit{pgfgantt}. 
En la plantilla pueden ver el código que lo genera y usarlo de base para construir el propio.

Las fechas pueden ser calculadas utilizando alguna de las herramientas antes citadas. Sin embargo, el siguiente ejemplo
fue elaborado utilizando 
\href{https://docs.google.com/spreadsheets/d/1fBz8NhSpc4tkkhz3KjJCbh1nR_ltDkfEcZi4tZXduqs}{esta hoja de cálculo}.

Es importante destacar que el ancho del diagrama estará dado por la longitud del texto utilizado para las tareas 
(Ejemplo: tarea 1, tarea 2, etcétera) y el valor \textit{x unit}. Para mejorar la apariencia del diagrama, es necesario
ajustar este valor y, quizás, acortar los nombres de las tareas.

\begin{figure}[htpb]
  \begin{center}
    \begin{ganttchart}[
      time slot unit=day,
      time slot format=isodate,
      x unit=0.038cm,
      y unit title=0.7cm,
      y unit chart=0.6cm,
      milestone/.append style={xscale=4}
      ]{2021-03-05}{2021-12-16}
      \gantttitlecalendar*{2021-03-05}{2021-12-16}{year} \\
      \gantttitlecalendar*{2021-03-05}{2021-12-16}{month} \\
      \ganttgroup{Duración Total}{2021-03-05}{2021-12-16} \\
      %%%%%%%%%%%%%%%%%Organización
      \ganttgroup{Organización}{2021-03-05}{2021-04-16} \\
      \ganttbar{Planificación del proyecto}{2021-03-05}{2021-04-15} \\
      %%%%%%%%%%%%%%%%%Ejecución
      \ganttgroup{Ejecución}{2021-04-16}{2021-10-21} \\
      \ganttbar{Tarea 1}{2021-04-16}{2021-04-29} \\
      \ganttbar{Tarea 2}{2021-04-30}{2021-05-13} \\
      \ganttbar{Tarea 3}{2021-05-14}{2021-05-27} \\
      \ganttbar{Tarea 4}{2021-05-28}{2021-07-12} \\
      \ganttbar{Tarea 5}{2021-07-13}{2021-08-09} \\
      \ganttbar{Tarea 6}{2021-08-10}{2021-09-23} \\
      \ganttbar{Tarea 7}{2021-09-24}{2021-09-30} \\
      \ganttbar{Tarea 8}{2021-10-01}{2021-10-14} \\
      \ganttbar{Tarea 9}{2021-10-15}{2021-10-21} \\
      % %%%%%%%%%%%%%%%%%Finalización
      \ganttgroup{Finalización}{2021-10-22}{2021-12-16} \\
      \ganttbar{Memoria v1}{2021-10-22}{2021-11-04} \\
      \ganttbar{Memoria v2}{2021-11-05}{2021-11-18} \\
      \ganttbar{Memoria final}{2021-11-19}{2021-12-02} \\
      % La fecha del siguiente milestone es la fecha en que terminamos la memoria
      \ganttmilestone{Enviar memoria al director}{2021-12-02} \\
      \ganttbar{Elaborar la presentación}{2021-12-03}{2021-12-16} \\
      \ganttmilestone{Ensayo de la presentación}{2021-12-16} \\
      %%%%%%%%%%%%%%%%%%%%%%%%%%%%%%%%%%%%%%%%%%%%%%%%%%%%%%%%%%%%%%%
    \end{ganttchart}
  \end{center}
  \caption{Diagrama de gantt de ejemplo}
  \label{fig:gantt}
\end{figure}


\begin{landscape}
\begin{figure}[htpb]
\centering 
\includegraphics[height=.85\textheight]{./Figuras/Gantt-2.png}
\caption{Ejemplo de diagrama de Gantt (apaisado).} %Modificar este título acorde.
\label{fig:diagGantt}
\end{figure}

\end{landscape}

\end{consigna}


\section{12. Presupuesto detallado del proyecto}
\label{sec:presupuesto}

\begin{consigna}{red}
Si el proyecto es complejo entonces separarlo en partes:
\begin{itemize}
	\item Un total global, indicando el subtotal acumulado por cada una de las áreas.
	\item El desglose detallado del subtotal de cada una de las áreas.
\end{itemize}

IMPORTANTE: No olvidarse de considerar los COSTOS INDIRECTOS.

Incluir la aclaración de si se emplea como moneda el peso argentino (ARS) o si se usa moneda extranjera (USD, EUR, etc). Si es en moneda extranjera se debe indicar la tasa de conversión respecto a la moneda local en una fecha dada.

\end{consigna}

\begin{table}[htpb]
\centering
\begin{tabularx}{\linewidth}{@{}|X|c|r|r|@{}}
\hline
\rowcolor[HTML]{C0C0C0} 
\multicolumn{4}{|c|}{\cellcolor[HTML]{C0C0C0}COSTOS DIRECTOS} \\ \hline
\rowcolor[HTML]{C0C0C0} 
Descripción &
  \multicolumn{1}{c|}{\cellcolor[HTML]{C0C0C0}Cantidad} &
  \multicolumn{1}{c|}{\cellcolor[HTML]{C0C0C0}Valor unitario} &
  \multicolumn{1}{c|}{\cellcolor[HTML]{C0C0C0}Valor total} \\ \hline
 &
  \multicolumn{1}{c|}{} &
  \multicolumn{1}{c|}{} &
  \multicolumn{1}{c|}{} \\ \hline
 &
  \multicolumn{1}{c|}{} &
  \multicolumn{1}{c|}{} &
  \multicolumn{1}{c|}{} \\ \hline
\multicolumn{1}{|l|}{} &
   &
   &
   \\ \hline
\multicolumn{1}{|l|}{} &
   &
   &
   \\ \hline
\multicolumn{3}{|c|}{SUBTOTAL} &
  \multicolumn{1}{c|}{} \\ \hline
\rowcolor[HTML]{C0C0C0} 
\multicolumn{4}{|c|}{\cellcolor[HTML]{C0C0C0}COSTOS INDIRECTOS} \\ \hline
\rowcolor[HTML]{C0C0C0} 
Descripción &
  \multicolumn{1}{c|}{\cellcolor[HTML]{C0C0C0}Cantidad} &
  \multicolumn{1}{c|}{\cellcolor[HTML]{C0C0C0}Valor unitario} &
  \multicolumn{1}{c|}{\cellcolor[HTML]{C0C0C0}Valor total} \\ \hline
\multicolumn{1}{|l|}{} &
   &
   &
   \\ \hline
\multicolumn{1}{|l|}{} &
   &
   &
   \\ \hline
\multicolumn{1}{|l|}{} &
   &
   &
   \\ \hline
\multicolumn{3}{|c|}{SUBTOTAL} &
  \multicolumn{1}{c|}{} \\ \hline
\rowcolor[HTML]{C0C0C0}
\multicolumn{3}{|c|}{TOTAL} &
   \\ \hline
\end{tabularx}%
\end{table}


\section{13. Gestión de riesgos}
\label{sec:riesgos}

\begin{consigna}{red}
a) Identificación de los riesgos (al menos cinco) y estimación de sus consecuencias:
 
Riesgo 1: detallar el riesgo (riesgo es algo que si ocurre altera los planes previstos de forma negativa)
\begin{itemize}
	\item Severidad (S): mientras más severo, más alto es el número (usar números del 1 al 10).\\
	Justificar el motivo por el cual se asigna determinado número de severidad (S).
	\item Probabilidad de ocurrencia (O): mientras más probable, más alto es el número (usar del 1 al 10).\\
	Justificar el motivo por el cual se asigna determinado número de (O). 
\end{itemize}   

Riesgo 2:
\begin{itemize}
	\item Severidad (S): X.\\
	Justificación...
	\item Ocurrencia (O): Y.\\
	Justificación...
\end{itemize}

Riesgo 3:
\begin{itemize}
	\item Severidad (S):  X.\\
	Justificación...
	\item Ocurrencia (O): Y.\\
	Justificación...
\end{itemize}


b) Tabla de gestión de riesgos:      (El RPN se calcula como RPN=SxO)

\begin{table}[htpb]
\centering
\begin{tabularx}{\linewidth}{@{}|X|c|c|c|c|c|c|@{}}
\hline
\rowcolor[HTML]{C0C0C0} 
Riesgo & S & O & RPN & S* & O* & RPN* \\ \hline
       &   &   &     &    &    &      \\ \hline
       &   &   &     &    &    &      \\ \hline
       &   &   &     &    &    &      \\ \hline
       &   &   &     &    &    &      \\ \hline
       &   &   &     &    &    &      \\ \hline
\end{tabularx}%
\end{table}

Criterio adoptado: 

Se tomarán medidas de mitigación en los riesgos cuyos números de RPN sean mayores a...

Nota: los valores marcados con (*) en la tabla corresponden luego de haber aplicado la mitigación.

c) Plan de mitigación de los riesgos que originalmente excedían el RPN máximo establecido:
 
Riesgo 1: plan de mitigación (si por el RPN fuera necesario elaborar un plan de mitigación).
  Nueva asignación de S y O, con su respectiva justificación:
  \begin{itemize}
	\item Severidad (S*): mientras más severo, más alto es el número (usar números del 1 al 10).
          Justificar el motivo por el cual se asigna determinado número de severidad (S).
	\item Probabilidad de ocurrencia (O*): mientras más probable, más alto es el número (usar del 1 al 10).
          Justificar el motivo por el cual se asigna determinado número de (O).
	\end{itemize}

Riesgo 2: plan de mitigación (si por el RPN fuera necesario elaborar un plan de mitigación).
 
Riesgo 3: plan de mitigación (si por el RPN fuera necesario elaborar un plan de mitigación).

\end{consigna}


\section{14. Gestión de la calidad}
\label{sec:calidad}

\begin{consigna}{red}
Elija al menos diez requerimientos que a su criterio sean los más importantes/críticos/que aportan más valor y para cada uno de ellos indique las acciones de verificación y validación que permitan asegurar su cumplimiento.

\begin{itemize} 
\item Req \#1: copiar acá el requerimiento con su correspondiente número.

\begin{itemize}
	\item Verificación para confirmar si se cumplió con lo requerido antes de mostrar el sistema al cliente. Detallar.
	\item Validación con el cliente para confirmar que está de acuerdo en que se cumplió con lo requerido. Detallar. 
\end{itemize}

\end{itemize}

Tener en cuenta que en este contexto se pueden mencionar simulaciones, cálculos, revisión de hojas de datos, consulta con expertos, mediciones, etc.  

Las acciones de verificación suelen considerar al entregable como ``caja blanca'', es decir se conoce en profundidad su funcionamiento interno.  

En cambio, las acciones de validación suelen considerar al entregable como ``caja negra'', es decir, que no se conocen los detalles de su funcionamiento interno.

\end{consigna}

\section{15. Procesos de cierre}    
\label{sec:cierre}

\begin{consigna}{red}
Establecer las pautas de trabajo para realizar una reunión final de evaluación del proyecto, tal que contemple las siguientes actividades:

\begin{itemize}
	\item Pautas de trabajo que se seguirán para analizar si se respetó el Plan de Proyecto original:\\
	 - Indicar quién se ocupará de hacer esto y cuál será el procedimiento a aplicar. 
	\item Identificación de las técnicas y procedimientos útiles e inútiles que se emplearon, los problemas que surgieron y cómo se solucionaron:\\
	 - Indicar quién se ocupará de hacer esto y cuál será el procedimiento para dejar registro.
	\item Indicar quién organizará el acto de agradecimiento a todos los interesados, y en especial al equipo de trabajo y colaboradores:\\
	  - Indicar esto y quién financiará los gastos correspondientes.
\end{itemize}

\end{consigna}

\end{document}